\section{直線でできることを曲線へ広げる}
長方形の面積は横の長さと縦の長さの積で求められる。直線の傾きも，横の変化量に対する縦の変化量の比で求められる。ところが，境界が曲線になったり，傾きが場所ごとに変わったりすると，同じ計算をそのまま使うことはできない。この本では，直線や長方形について知っている計算から出発し，曲線の変化と面積を調べる方法を作る。

まず，二つの問いを並べてみよう。第一は，放物線 $y=x^2$ 上の点 $(1,1)$ で，曲線はどの程度の傾きで進んでいるのかという問いである。第二は，この放物線と $x$ 軸，直線 $x=1$ に囲まれた部分の面積はいくらかという問いである。第一の問いでは二点を結ぶ直線を考え，その二点を近づける。第二の問いでは細い長方形を並べ，その幅を小さくする。どちらも，計算できるものから始めて，細かくした先にある値を求めようとしている。

ここで「限りなく近づける」という表現は，まだ計算の見通しを説明するために使っている。何をどのように近づけるのかを明らかにするのが\chapterref{章:極限と微分}であり，長方形の和を積分として定義するのが\chapterref{章:定積分}である。この章では具体的な図と計算を通して，後の定義が何のために必要なのかを確かめる。

